\documentclass[12pt,a4paper,sans]{moderncv}
\moderncvstyle{classic}
\moderncvcolor{blue}

\usepackage[utf8]{inputenc}
% \usepackage[scale=0.90]{geometry}
\usepackage[top=1.5cm,  bottom=1.2cm, left=1.5cm, right=1.5cm]{geometry}
\usepackage{multicol}

\definecolor{navyblue}{RGB}{0, 0, 128}
\usepackage[colorlinks=true, urlcolor=navyblue]{hyperref}

\colorlet{bodyrulecolor}{blue}
\colorlet{sectioncolor}{blue}

\usepackage{enumitem}  % for item sizes
\setlist[itemize]{topsep=4pt, itemsep=3pt, parsep=0pt}   % for item sizes


\name{Javad}{Komijani}
% \title{Senior Computational Physicist}
\title{Computational Scientist \& ML Researcher}
\address{Zurich, Switzerland}
\phone[mobile]{+41 76 796 9992}
\email{jkomijani@ethz.ch}
\homepage{jkomijani.github.io}
\social[github]{jkomijani}
% \social[linkedin]{javad-komijani-543a31235}


\begin{document}
\makecvtitle


\section{Professional Summary}

Computational scientist and machine learning researcher specializing in generative models, diffusion processes, and normalizing flows on structured and symmetry-constrained domains. Experienced in building PyTorch-based scientific ML frameworks, distributed training systems, and GPU-accelerated simulation pipelines. Focus on scalable generative modeling, geometric deep learning, and scientific computing at scale.

% Senior computational physicist with over 10 years of experience in scientific computing, 5 years in scientific machine learning, 
% and an extensive background in mathematical methods, numerical simulations, theoretical modeling, statistical analysis.
% Holds an MSc in Electrical Engineering and a PhD in Physics.
% Develops generative models (diffusion models, normalizing flows) in PyTorch for Lie groups. 
% Passionate about applying data science and machine learning to solve complex, real-world problems.
% Expert in differential equations, statistical analysis, and numerical optimization. 
% Bridges advanced mathematical modeling with machine learning to solve complex computational problems.

 
% ==========================
\section*{Selected Research Impact}

\begin{itemize}[leftmargin=*]
    \item 26 peer-reviewed publications, 3000+ citations, h-index 19.
    \item Contributions to precision Standard Model parameter determinations (FLAG / PDG collaborations).
    \item Developed frameworks for controlling multi-scale expansions in quantum field theory.
    \item Introduced methods for handling divergent perturbative series in physics computations.
    \item Developed an implicit score-matching framework for Lie groups.
    \item Developed GPU-accelerated software frameworks for diffusion models and normalizing flows.
\end{itemize}


\section*{Research \& Technical Experience}

\cventry{2021--2026}{Postdoctoral Researcher}{ETH Zurich}{}{}{
\begin{itemize}
\item Developed diffusion and score-based generative models on Lie groups for lattice gauge theory.
\item Built differentiable computing frameworks using PyTorch (ODE solvers and linear algebra).
\item Integrated normalizing flow modules into the EU interTwin Digital Twin Engine.
\item Developed distributed training infrastructure for large-scale lattice simulations.
\item Supervised PhD, MSc, BSc, and internship projects in scientific machine learning.
\end{itemize}
}

\cventry{2015--2021}{Postdoctoral Researcher}{TU Munich / Glasgow / Tehran}{}{}{
\vspace{-0.5em}
\begin{itemize}
\item Monte Carlo simulations for lattice QCD observable.
\item Bayesian parameter estimation from Lattice QCD simulations.
% \item Developed models and applied Bayesian inference \& Monte Carlo simulations to large-scale datasets.
\item Developed scalable lattice-QCD simulation pipelines in HPC environments.
\item Collaborated internationally in computational and mathematical physics projects.
\end{itemize}
}


%--------------------------------------------------
\section{Open-Source Scientific ML Software}

\cventry{\href{https://github.com/jkomijani/lattice_ml}{lattice\_ml}}{Installable Python package for scientific machine learning}{}{}{}{
\begin{itemize}
\item Diffusion-based generative modeling on Lie groups and gauge theories.
\item Adjoint-based differentiable ODE solvers for scalar and Lie-group-valued systems.
\item Differentiable SVD and eigendecomposition of unitary matrices.
\end{itemize}
}
\cvitem{\href{https://github.com/jkomijani/normflow_}{normflow}}{Installable Python package for normalizing flows on Lie groups and lattice gauge theory, including symmetry-preserving transformations and flow-based sampling methods.}
\cvitem{\href{https://www.intertwin.eu/article/thematic-module-normflow}{interTwin}}{Contributed normalizing-flow-based ML modules to the EU interTwin Digital Twin Engine.}


\clearpage

% ==========================
\section*{Mentoring \& Leadership}

\cventry{Mentoring}{Co-supervised PhD, MSc, and BSc researchers in physics \& scientific ML:}{}{}{}{
     \begin{itemize}
    \item Gauge-equivariant diffusion models for lattice field theory (PhD project)
    \item Fourier acceleration methods for SU(N) lattice gauge theory (PhD project)
    \item Autoregressive generative models for structured lattice configurations (MSc project)
    \item Continuous-time generative models and normalizing flows (MSc projects)
    \item Stability and sampling improvements for topological systems (MSc project)
    \item CNN architectures in arbitrary dimensions for scientific data (Intern project)
\end{itemize}
}

\cventry{Organization}{Co-organized workshops on physics \& machine learning:}{}{}{}{
\begin{itemize}
\item ZPW 2026 — Machine Learning Techniques for Particle Physics.
\item QCD 2025 — Symposium on QCD Factorization.
\item ML Lattice 2025 — Machine learning approaches in Lattice QCD.
\end{itemize}
}

%--------------------------------------------------
\section{Technical Skills}

\cvitem{Programming}{Python, C++, Cython, Bash}
\cvitem{ML/AI}{PyTorch, Diffusion Models, Normalizing Flows, Generative Modeling, Distributed Training}
\cvitem{Mathematics}{Differential Equations, Linear Algebra, Optimization, Statistical Modeling}
\cvitem{Tools}{Linux, Git, SQL, Pandas}
\cvitem{Scientific}{NumPy, SciPy, Scikit-learn, HPC, Parallel Computing, Monte Carlo Methods}


% =================
\section{Education}

\cventry{2015}{PhD in Physics}{Washington University in St. Louis}{}{}{
\textbf{Thesis}: \textit{Topics in Lattice Gauge Theory and Theoretical Physics}
\begin{itemize}
  \item Developed EFT-based frameworks and Bayesian inference methods for lattice QCD.
  \item Performed precision determinations of hadronic observables in lattice QCD.
  \item Developed a novel definition of eigenvalue problems for nonlinear systems arising in physics.
\end{itemize}
}
\cventry{2009}{MSc in Electrical Engineering}{University of Tehran}{}{}{
\begin{itemize}
\item Focus: electromagnetic waves and fields, signal processing, statistical analysis.
\end{itemize}
}
\cventry{2006}{BSc in Electrical Engineering}{University of Tehran}{}{}{}

%--------------------------------------------------
\section{Publications}

\cvitem{}{A complete list of over 26 peer-reviewed publications can be found at \href{https://jkomijani.github.io/docs/Publications.pdf}{this link}.}

%--------------------------------------------------
\section{Languages}

\cvitem{English}{Fluent}
\cvitem{German}{Basic (A2)}


\clearpage

\end{document}